\section{The separation theorem}
\label{sec:synthesis}

\begin{proof}[Proof of Theorem~\ref{thm:separation-family}]
Corollary~\ref{cor:universal-ch0} gives universal \(CH_0\)-triviality of
\(X_b\).  The projective bundle formula for Chow groups gives the same
property for \(X_b\times\PP^1\).  Theorem~\ref{thm:every-cubic} proves that
this fourfold is irrational.  Non-isotriviality follows from the nonconstant
intermediate-Jacobian period map on the \(A_5\)-component.
\end{proof}

The theorem separates two tests on one variety.  The algebraic class
\(\Theta^4/4!\) removes the decomposition-of-diagonal obstruction.  The
primitive-sixth framed-monodromy multiplicity survives one stabilization and
obstructs a rational factorization.  The first mechanism is special to the
exotic \(A_5\)-family;
the second holds for every cubic threefold.  More generally, the local
cofactor-saturation theorem applies beyond this family; this geometry is where
the two mechanisms hold simultaneously.

The result stops after one stabilization.  For
\(X\times\PP^m\) with \(m\ge2\), weak-factorization centers may have dimension
at least three and can carry primitive-sixth quantum packets.  The
low-dimensional vanishing in
Proposition~\ref{prop:low-dimensional-vanishing} no longer proves
noncancellation.
Likewise, the fibrewise minimal-class argument gives no relative universal
cycle over the pencil.  These are separate problems and are not implicit in
the theorem above.

Thus the same smooth fourfolds pass the universal \(CH_0\)-triviality test
and fail the framed-quantum rationality test.  Vanishing of the
decomposition-of-the-diagonal obstruction does not predict rationality even
after one explicit stabilization.
